\section{Adaptive recovery of complete component boundaries}
\label{sec:boundary}
We first acquire a coarse, fixed-resolution description. Its purpose
is to find interior centers and isolated radial brackets, not to supply
the final accuracy. This avoids both a supplied free reference point
and a bisection that might accidentally encounter a different obstacle.

\begin{lemma}[Coarse hulls and interior centers]\label{lem:coarse}
A prior-dependent finite number of queries from
Proposition~\ref{prop:query}, at one fixed scale $\sigma_c>0$, recovers
all complete components needed in the protected patch to Hausdorff
error $e_c\le3\sigma_c$. It supplies, for each such component $C$, a
polygon $H$ and a computable center $o$ for which
\begin{equation}\label{eq:inner-ball}
 \overline B_r(o)\subset C\cap H\subset\overline B_{R_1}(o)
\end{equation}
with fixed positive $r,R_1$. For every direction $v$, it supplies a
radial interval of fixed length at most $C e_c$ containing the boundary
point of $C$ on $o+\R_+v$. The interval is disjoint from the
$\sigma$-enlargements of all other components for every sufficiently
small fine scale $\sigma$.
\end{lemma}
\begin{proof}
The curvature upper bound supplies an interior rolling radius
$r_* =1/\kappa_+$. In support coordinates this follows from
$p+p''\ge r_*$: subtracting the support function of the radius-$r_*$
disk leaves a support function with nonnegative curvature measure.
In particular $C=C_{-2\sigma_c}+2\sigma_c\overline B$ when
$2\sigma_c<r_*$. Choose $\sigma_c<\min(r_*/8,d_0/40)$ and a square
grid of covering radius $\rho_c\le\sigma_c/8$ on a surrounding
square with collar larger than $2D_0+2$.

At its nodes retain the estimated occupations at least $1/2$, join
retained nodes at distance at most $4\sigma_c$, and take a convex hull
for each connected cluster. Every retained node is within $\sigma_c$
of a unique body; all nodes of $C_{-\sigma_c}$ are retained. For any
retained node assigned to $C$, choose a point in $C_{-2\sigma_c}$ at
distance at most $3\sigma_c$, and then a nearest grid node. That node
is retained and joined to the first. Convexity connects any two such
deep nodes by a segment; partitioning it into lengths at most $\rho_c$
and replacing partition points by nearest nodes yields a retained
path with successive lengths at most $3\rho_c$. Different bodies
cannot join because $d_0-2\sigma_c>4\sigma_c$. The hull thus lies in
$C+\sigma_c\overline B$, whereas every point of $C$ is within
$2\sigma_c+\rho_c$ of a retained deep node. This proves the error.
Discard clusters whose hulls are within $D_0+1$ of the outer square
boundary. Fragments cannot survive this cutoff, and the chosen collar
protects every required complete component.

Compute a largest inscribed disk of $H$ by its finite half-plane
inequalities, to any fixed reserved precision. Since
$d_H(H,C)\le e_c$, the support inequalities imply
$C_{-e_c}\subset H$ and $H_{-e_c}\subset C$. A disk of radius $r_*$
in $C$ therefore gives an inscribed disk of radius at least $r_*-e_c$
in $H$. The computed center, with the radii reduced by the error and
numerical reserve, satisfies \eqref{eq:inner-ball} for a fixed $r>0$.
Diameter bounds supply $R_1$.

Write $R_H(v),R_C(v)$ for the radial functions about $o$. Since both
bodies contain $\overline B_r(o)$,
\[
 H+e_c\overline B\subset o+(1+e_c/r)(H-o),
\]
and the analogous inclusion holds for $C$. The two Hausdorff
inclusions give $|R_H(v)-R_C(v)|\le C e_c$ uniformly. The polygonal
radial value is computable from its sides. Use a slightly enlarged
bracket about it, allowing numerical error. By decreasing the fixed
$\sigma_c$, its length and every point's distance to $C$ are less
than $d_0/4$. Separation excludes every other component from the
bracket and its sufficiently small kernel neighborhood.
\end{proof}

\begin{lemma}[Bisection with an indeterminate boundary layer]
\label{lem:bisection}
On each bracket in Lemma~\ref{lem:coarse}, the derived label of
Proposition~\ref{prop:query} may be arbitrary, and may depend on the
past, only when
\begin{equation}\label{eq:radial-ambiguity}
                       |a-R_C(v)|\le A\sigma
\end{equation}
for a fixed $A$. After $k$ midpoint tests, ordinary bisection returns
an estimate with error at most $A\sigma+2^{-k-1}w_0$, where $w_0$
is the initial bracket length. In particular $O(\log(C/\sigma))$
queries recover $R_C(v)$ to $O(\sigma)$.
\end{lemma}
\begin{proof}
Translate $o$ to zero. The inner ball implies
\[
 (1-\sigma/r)C\subset C_{-\sigma},\qquad
 C+\sigma\overline B\subset(1+\sigma/r)C.
\]
Together with $R_C\le R_1$, these show that outside
\eqref{eq:radial-ambiguity} the label is the true radial inside/outside
label. Other components were excluded from the bracket. A label one
at radius $a$ implies $a\le R_C+A\sigma$; a label zero implies
$a\ge R_C-A\sigma$. If the current endpoints are $l,u$, midpoint
updates consequently preserve the relaxed invariant
$R_C\in[l-A\sigma,u+A\sigma]$. This does not require $R_C$ to remain
in $[l,u]$, or the noisy labels to be monotone. The width halves at
every step. Its midpoint has the asserted error.
\end{proof}

\begin{lemma}[Radial regularity and stable interpolation]
\label{lem:radial}
Under \eqref{eq:inner-ball} and the physical support priors, the
$2\pi$-periodic radial function $R_C(\varphi)$ has a uniformly bounded
$C^{6,\beta}$ norm. Let $h=2\pi/m$ for an integer $m\ge9$. From
radial values at the $m$ equally spaced angles, each with error at most
$e$, one can construct a smooth function $\widehat R$ such that
\begin{equation}\label{eq:radial-interpolation}
 \|\widehat R-R_C\|_{C^j}
           \le C(h^{s-j}+e h^{-j}),\qquad 0\le j\le3.
\end{equation}
For $e\le C_0h^s$ and sufficiently small $h$, the curve
$o+\widehat R(\varphi)(\cos\varphi,\sin\varphi)$ is the boundary of
a smooth strictly convex body $\widehat C$, and
\begin{equation}\label{eq:geometric-interpolation}
 d_H(\widehat C,C)\le Ch^s,\qquad
 \|p_{\widehat C}^{\rm lab}-p_C^{\rm lab}\|_{C^2}\le Ch^{s-2}.
\end{equation}
Here uncentered support functions are compared in the same laboratory
frame; centering preserves these orders.
\end{lemma}
\begin{proof}
Use the support function $p$ relative to $o$, so $p\ge r$. In normal
angle $\theta$, the boundary is $p n+p' n^\perp$, its radial angle is
\[
 \varphi=\theta+\arctan(p'/p),\qquad
 \frac{\dd\varphi}{\dd\theta}
             =\frac{p(p+p'')}{p^2+(p')^2}>c>0.
\]
The inverse $\theta(\varphi)$ is uniformly $C^{5,\beta}$. Although
the expression $R=\sqrt{p^2+(p')^2}$ initially gives only five
ordinary derivatives, differentiation in radial angle cancels the
curvature factor:
\begin{equation}\label{eq:radial-cancellation}
                   R'=R(p'/p)\circ\theta.
\end{equation}
Its right side is uniformly $C^{5,\beta}$. Thus $R$ is uniformly
$C^{6,\beta}$, without losing a derivative from the support prior.

At each angular node interpolate the seven consecutive sampled
values by a polynomial of degree at most six in a local lift of the
angle. Combine these polynomials with a fixed translated smooth
partition of unity supported on neighborhoods of size $O(h)$.
On each such neighborhood the normalized Vandermonde inverse is
bounded. A degree-six Taylor polynomial of the true function is
reproduced exactly, and its remainder and derivatives through order
three are $O(h^{s-j})$. Errors in the values contribute $O(eh^{-j})$.
Derivatives falling on the partition have size $O(h^{-j})$ and are
absorbed by the corresponding lower-order remainders. The overlap
number is fixed. This proves \eqref{eq:radial-interpolation}, including
across the angular cut by periodic lifts.

A positive radial function is strictly convex when
\[
       F=R^2+2(R')^2-RR''>0.
\]
For the true curve this expression has a uniform positive margin,
by the inradius, curvature and diameter bounds. It persists for
$\widehat R$ by the $C^2$ estimate. The two radial graphs are uniformly
$O(h^s)$-close, which proves the Hausdorff bound.

To check the stronger support norm, set
$\theta=\varphi-\arctan(R'/R)$, now viewing the normal angle as a
function of radial angle. Its derivative is
$F/(R^2+(R')^2)>c$. At corresponding angles the support value and
radius of curvature are
\[
 p=\frac{R^2}{\sqrt{R^2+(R')^2}},\qquad
 p+p''=\frac{(R^2+(R')^2)^{3/2}}{F}.
\]
These expressions, and the tangential support derivative, depend
Lipschitzly on $R,R',R''$ on the bounded set with these margins.
The inverse normal-angle maps differ by $O(\|\widehat R-R\|_{C^1})$.
Their composition changes the true expressions by the same order,
using the uniform $C^3$ bound; the reconstructed functions also have
such a bound by \eqref{eq:radial-interpolation}. Thus the support
error through two derivatives is $O(h^{s-2})$. Translation to the
laboratory frame adds the same first harmonic to both supports. The
Steiner point is a bounded linear functional of the support, so its
error is $O(h^s)$ and centering preserves the conclusions.
\end{proof}

\begin{proposition}[Adaptive patch reconstruction]\label{prop:patch}
For a fixed protected patch, complete component boundaries can be
recovered with the errors in \eqref{eq:geometric-interpolation} using
at most $C h^{-1}\log(C/h)$ local occupation queries. The finest
localization is $\sigma=c h^s$. With probability at least $1-\delta$
the attempted-bit cost is at most
\begin{equation}\label{eq:patch-cost}
           C h^{-1}\log(C/h)\log(C/(h\delta)).
\end{equation}
The number of components, the initial coarse queries and every
local Bellman stencil size are bounded solely by the fixed prior.
\end{proposition}
\begin{proof}
On the simultaneous accuracy event, obtain the hulls and centers in
Lemma~\ref{lem:coarse}. Separation and the fixed aperture bound their
number. For each body use $m$ equally spaced radial directions and
their isolated coarse brackets. Lemma~\ref{lem:bisection} uses
$O(\log(C/h))$ midpoint queries per direction to return values with
error $e\le C\sigma$. Lemma~\ref{lem:radial} gives the claimed
geometric estimates. The total number of queries is bounded by a
deterministic $Q_h=C(1+h^{-1}\log(C/h))$. Allocate failure probability
$\delta/Q_h$ to every query, including the coarse ones, before the
experiment. Conditional confidence and the union bound in
Proposition~\ref{prop:query} give the simultaneous event. Its
preparation cost is \eqref{eq:patch-cost}. Locations chosen using
coarse or previous fine data do not alter this argument.

For a finite numerical realization, polygon intersections, midpoint
locations and radial values are enclosed with error $O(\sigma)$.
The estimate \eqref{eq:radial-interpolation} absorbs those value errors.
Represent the output by its finite polynomial coefficients and the
fixed smooth partition functions. They, their derivatives and the
normal-angle maps can be evaluated to any positive prescribed error
by finite operations on compact intervals with the stated derivative
bounds. This constructs an actual smooth convex output, without
assigning a false smoothness bound to the initial polygons. No
infinite-dimensional candidate minimization is performed.
\end{proof}
